\Article{29}
Before going further, we should note that the criteria for prime numbers (taken positively when they have the form $4n+1$, negatively when they have the form $4n+3$) suffice to decide the question for all other numbers, provided that the theorem of Article 7 and the criteria for $-1$ and $\pm2$ are also used. Thus, for example, if criteria for the number $+3$ are wanted, the criteria given in Article 25, which concern $-3$, will also hold for $+3$ whenever $\frac12b$ is even. The collections $A,B,C,D$ must instead be exchanged with the collections $C,D,A,B$ whenever $\frac12b$ is odd, giving the following rules:


\SourceLeaf{111}%
\begin{center}$+3$ belongs\end{center}
{\small
\[
\begin{array}{c|l}
\substack{\text{to the}\\\text{collection}}&\text{if}\\
A&\begin{aligned}b\equiv0\ (\mathrm{mod.}\ 12);\\\text{or simultaneously }a\equiv0\ (\mathrm{mod.}\ 3),\ b\equiv2\ (\mathrm{mod.}\ 4)\end{aligned}\\
B&b\equiv8a\text{ or }10a\ (\mathrm{mod.}\ 12)\\
C&\begin{aligned}b\equiv6a\ (\mathrm{mod.}\ 12);\\\text{or simultaneously }a\equiv0\ (\mathrm{mod.}\ 3),\ b\equiv0\ (\mathrm{mod.}\ 4)\end{aligned}\\
D&b\equiv2a\text{ or }4a\ (\mathrm{mod.}\ 12)
\end{array}
\]
}

Likewise, the criteria for $\pm6$ are obtained by combining the criteria for $\mp2$ and $-3$; namely,
\begin{center}$+6$ belongs\end{center}
{\small
\[
\begin{array}{c|l}
\substack{\text{to the}\\\text{collection}}&\text{if}\\
A&\begin{aligned}b\equiv0,\quad2a,22a\ (\mathrm{mod.}\ 24);\\\text{or simultaneously }a\equiv0\ (\mathrm{mod.}\ 3),\ b\equiv4a\ (\mathrm{mod.}\ 8)\end{aligned}\\
B&\begin{aligned}b\equiv4a,\quad6a,\quad8a\ (\mathrm{mod.}\ 24);\\\text{or simultaneously }a\equiv0\ (\mathrm{mod.}\ 3),\ b\equiv2a\ (\mathrm{mod.}\ 8)\end{aligned}\\
C&\begin{aligned}b\equiv10a,12a,14a\ (\mathrm{mod.}\ 24);\\\text{or simultaneously }a\equiv0\ (\mathrm{mod.}\ 3),\ b\equiv0\ (\mathrm{mod.}\ 8)\end{aligned}\\
D&\begin{aligned}b\equiv16a,18a,20a\ (\mathrm{mod.}\ 24);\\\text{or simultaneously }a\equiv0\ (\mathrm{mod.}\ 3),\ b\equiv6a\ (\mathrm{mod.}\ 8)\end{aligned}
\end{array}
\]
}

\begin{center}while $-6$ belongs\end{center}
{\small
\[
\begin{array}{c|l}
\substack{\text{to the}\\\text{collection}}&\text{if}\\
A&\begin{aligned}b\equiv0,\quad10a,14a\ (\mathrm{mod.}\ 24);\\\text{or simultaneously }a\equiv0\ (\mathrm{mod.}\ 3),\ b\equiv4a\ (\mathrm{mod.}\ 8)\end{aligned}\\
B&\begin{aligned}b\equiv4a,\quad8a,18a\ (\mathrm{mod.}\ 24);\\\text{or simultaneously }a\equiv0\ (\mathrm{mod.}\ 3),\ b\equiv6a\ (\mathrm{mod.}\ 8)\end{aligned}\\
C&\begin{aligned}b\equiv2a,12a,22a\ (\mathrm{mod.}\ 24);\\\text{or simultaneously }a\equiv0\ (\mathrm{mod.}\ 3),\ b\equiv0\ (\mathrm{mod.}\ 8)\end{aligned}\\
D&\begin{aligned}b\equiv6a,16a,20a\ (\mathrm{mod.}\ 24);\\\text{or simultaneously }a\equiv0\ (\mathrm{mod.}\ 3),\ b\equiv2a\ (\mathrm{mod.}\ 8)\end{aligned}
\end{array}
\]
}

In the same way, the criteria for the number $+21$ will be assembled from the criteria for $-3$ and $-7$; the criteria for $-105$ from those for $-1$, $-3$, $+5$, $-7$, etc.

\Article{30}
Inductive investigation thus opens up a very rich harvest of particular theorems related to the theorem for the number $2$. But a common bond is lacking, rigorous proofs are lacking, since the method by which we completed the treatment of the number $2$ in the first paper admits no further application. There are indeed various methods by which proofs could be obtained for particular cases, especially those concerning the distribution of quadratic residues between the collections $A,C$. We do not dwell on them, however, since a general\ \SourceLeaf{112}%
 theory embracing \emph{all} cases must be our aim. Having begun to devote our thoughts to this matter in $1805$, we soon became convinced that the true source of the general theory must be sought in an extension of the domain of arithmetic, as we already indicated in Article 1.

While higher arithmetic, in the questions treated so far, deals only with real integers, the theorems on biquadratic residues shine in their full simplicity and genuine beauty only when the domain of arithmetic is extended to \emph{imaginary} quantities. Its objects are then, without restriction, numbers of the form $a+bi$, where $i$ denotes, as usual, the imaginary quantity $\surd-1$, and $a,b$ range without restriction over all real integers between $-\infty$ and $+\infty$. We shall call such numbers \emph{complex integers}, treating real numbers as a species contained within the complex numbers rather than as opposed to them. The present paper will set out both the elementary theory of complex numbers and the first beginnings of the theory of biquadratic residues, which we shall undertake to complete in every respect in a subsequent continuation\begingroup\renewcommand{\thefootnote}{*)}\footnote{We should at least mention here in passing that the domain thus defined is particularly suited to the theory of biquadratic residues. The theory of cubic residues must similarly be built on the consideration of numbers of the form $a+bh$, where $h$ is an imaginary root of the equation $h^3-1=0$, say $h=-\frac12+\surd\frac34\pdot i$; and likewise the theory of residues of higher powers will require the introduction of other imaginary quantities.}\endgroup.

\clearpage
\Article{31}
First we introduce some names that will help us to be brief and clear.

The domain of complex numbers $a+bi$ contains
\begin{list}{}{\setlength{\leftmargin}{1.6em}\setlength{\labelwidth}{1.4em}\setlength{\labelsep}{.2em}}
\item[I.] real numbers, where $b=0$, and among these, according to the nature of $a$,
\begin{list}{}{\setlength{\leftmargin}{1.7em}\setlength{\labelwidth}{1.5em}\setlength{\labelsep}{.2em}}
\item[1)] zero,
\item[2)] positive numbers,
\item[3)] negative numbers;
\end{list}
\item[II.] imaginary numbers, where $b$ is unequal to zero. Here we distinguish in turn
\begin{list}{}{\setlength{\leftmargin}{1.7em}\setlength{\labelwidth}{1.5em}\setlength{\labelsep}{.2em}}
\item[1)] imaginary numbers without a real part, that is, where $a=0$;
\item[2)] imaginary numbers with a real part, where $b$ is nonzero and $a=0$ also does not hold.
\end{list}
\end{list}
The former may, if desired, be called pure imaginary numbers, and the latter mixed imaginary numbers.

\SourceLeaf{113}%
In this theory we use four units, $+1$, $-1$, $+i$, $-i$, which will be called simply positive, negative, positive imaginary and negative imaginary.

We shall call the three products of any complex number by $-1$, $+i$, $-i$ its \emph{companions} or \emph{numbers associated with it}. Thus, except for zero (which is associated with itself), four \emph{distinct} numbers are always associated.

By contrast, we call the number obtained from a complex number by exchanging $i$ with $-i$ its \emph{conjugate}. Among imaginary numbers, therefore, two \emph{distinct} numbers are always conjugate, while real numbers are their own conjugates, if we choose to extend the name to them.

We call the product of a complex number and its conjugate the \emph{norm} of each. The square of a real number is therefore to be taken as its norm.

In general we have eight related numbers, namely,
\[
\begin{array}{r|r}
a+bi&a-bi\\
-b+ai&-b-ai\\
-a-bi&-a+bi\\
b-ai&b+ai
\end{array}
\]
in which we see two groups of four associated numbers and four pairs of conjugates; their common norm is $aa+bb$. But the eight numbers reduce to four distinct ones whenever either $a=\pm b$, or $a,b=0$ for one of the two numbers.

\Needspace{5\baselineskip}
The following results follow immediately from these definitions:

The conjugate of the product of two complex numbers is the product of their conjugates.

The same holds for a product of more factors, and also for quotients.

The norm of the product of two complex numbers is equal to the product of their norms.

This theorem also extends to products of any number of factors and to quotients.

The norm of any complex number (except zero, an exception that from now on we shall generally leave understood) is a \emph{positive} number.

Moreover, nothing prevents us from extending our definitions to fractional or even irrational values of $a,b$; but $a+bi$ will be called a complex integer only when \emph{each} of $a,b$ is an integer, and rational only when \emph{each} of $a,b$ is rational.

